\begin{ex}
Tính kết quả của phép tính: $47 - 5$.
\shortans{42}
\loigiai{
Hàng đơn vị: $7 - 5 = 2$. Giữ nguyên 4 chục: $47 - 5 = 42$.
\textbf{Đáp số:} $42$.
}
\end{ex}

\begin{ex}
Tính kết quả của phép tính: $68 - 6$.
\shortans{62}
\loigiai{
Hàng đơn vị: $8 - 6 = 2$. Giữ nguyên 6 chục: $68 - 6 = 62$.
\textbf{Đáp số:} $62$.
}
\end{ex}

\begin{ex}
Tính kết quả của phép tính: $89 - 4$.
\shortans{85}
\loigiai{
Hàng đơn vị: $9 - 4 = 5$. Giữ nguyên 8 chục: $89 - 4 = 85$.
\textbf{Đáp số:} $85$.
}
\end{ex}

\begin{ex}
Điền số thích hợp vào ô trống: $56 - \dots = 51$.
\shortans{5}
\loigiai{
Muốn tìm số trừ, lấy: $56 - 51 = 5$.
\textbf{Đáp số:} $5$.
}
\end{ex}

\begin{ex}
Điền số thích hợp vào ô trống: $\dots - 3 = 74$.
\shortans{77}
\loigiai{
Muốn tìm số bị trừ, lấy: $74 + 3 = 77$.
\textbf{Đáp số:} $77$.
}
\end{ex}

\begin{ex}
Tính giá trị của biểu thức: $78 - 3 - 2$.
\shortans{73}
\loigiai{
Thực hiện: $78 - 3 = 75$; $75 - 2 = 73$.
\textbf{Đáp số:} $73$.
}
\end{ex}

\begin{ex}
Tính giá trị của biểu thức: $96 - 4 - 2$.
\shortans{90}
\loigiai{
Thực hiện: $96 - 4 = 92$; $92 - 2 = 90$.
\textbf{Đáp số:} $90$.
}
\end{ex}

\begin{ex}
Lớp 1B có 38 học sinh, trong đó có 5 học sinh vắng mặt. Hỏi lớp 1B có bao nhiêu học sinh đang có mặt ở lớp?
\shortans{33}
\loigiai{
Số học sinh có mặt là: $38 - 5 = 33$ (học sinh).
\textbf{Đáp số:} $33$.
}
\end{ex}

\begin{ex}
Một cây dừa có 29 quả bắp dừa. Bác nông dân hái bớt 7 quả. Hỏi trên cây dừa còn lại bao nhiêu quả?
\shortans{22}
\loigiai{
Số quả dừa còn lại là: $29 - 7 = 22$ (quả).
\textbf{Đáp số:} $22$.
}
\end{ex}

\begin{ex}
Một cửa hàng có 68 chiếc xe máy, đã bán được 5 chiếc xe máy. Hỏi cửa hàng còn lại bao nhiêu chiếc xe máy?
\shortans{63}
\loigiai{
Số xe máy còn lại là: $68 - 5 = 63$ (chiếc xe máy).
\textbf{Đáp số:} $63$.
}
\end{ex}

\begin{ex}
Tìm số tự nhiên $x$ thỏa mãn: $87 - x = 82$.
\shortans{5}
\loigiai{
Ta có: $x = 87 - 82 = 5$.
\textbf{Đáp số:} $5$.
}
\end{ex}

\begin{ex}
Tìm số tự nhiên $y$ thỏa mãn: $y - 7 = 41$.
\shortans{48}
\loigiai{
Ta có: $y = 41 + 7 = 48$.
\textbf{Đáp số:} $48$.
}
\end{ex}

\begin{ex}
Một sợi dây dài 59 cm. Bạn Linh cắt bớt một đoạn dài 6 cm. Hỏi sợi dây còn lại dài bao nhiêu xăng-ti-mét?
\shortans{53}
\loigiai{
Chiều dài sợi dây còn lại là: $59 - 6 = 53$ (cm).
\textbf{Đáp số:} $53$.
}
\end{ex}

\begin{ex}
Quyển truyện có 76 trang. Bạn Nam đã đọc được 4 trang. Hỏi Nam còn phải đọc bao nhiêu trang nữa?
\shortans{72}
\loigiai{
Số trang Nam còn phải đọc là: $76 - 4 = 72$ (trang).
\textbf{Đáp số:} $72$.
}
\end{ex}

\begin{ex}
Tìm một số biết rằng lấy 88 trừ đi số đó thì được 81.
\shortans{7}
\loigiai{
Số cần tìm là: $88 - 81 = 7$.
\textbf{Đáp số:} $7$.
}
\end{ex}

\begin{ex}
Hai số có hiệu bằng 4, số lớn là 69. Hỏi số bé là số mấy?
\shortans{65}
\loigiai{
Số bé là: $69 - 4 = 65$.
\textbf{Đáp số:} $65$.
}
\end{ex}

\begin{ex}
Số 58 trừ đi số nào để được 52?
\shortans{6}
\loigiai{
Số đó là: $58 - 52 = 6$.
\textbf{Đáp số:} $6$.
}
\end{ex}

\begin{ex}
Tính: $49 - 2 - 4$.
\shortans{43}
\loigiai{
Thực hiện: $49 - 2 = 47$; $47 - 4 = 43$.
\textbf{Đáp số:} $43$.
}
\end{ex}

\begin{ex}
Vườn nhà bà có 48 cây bắp cải, bà đã thu hoạch 6 cây. Hỏi trong vườn còn lại bao nhiêu cây bắp cải?
\shortans{42}
\loigiai{
Số cây bắp cải còn lại là: $48 - 6 = 42$ (cây).
\textbf{Đáp số:} $42$.
}
\end{ex}

\begin{ex}
Một đàn gà có 39 con gà. Có 7 con gà chạy ra vườn kiếm ăn. Hỏi trong chuồng còn lại bao nhiêu con gà?
\shortans{32}
\loigiai{
Số con gà trong chuồng còn lại là: $39 - 7 = 32$ (con).
\textbf{Đáp số:} $32$.
}
\end{ex}

\begin{ex}
Tìm số tự nhiên $a$ sao cho: $a + 4 = 79 - 1$.
\shortans{74}
\loigiai{
Vế phải: $79 - 1 = 78$. Ta có: $a + 4 = 78 \Rightarrow a = 78 - 4 = 74$.
\textbf{Đáp số:} $74$.
}
\end{ex}

\begin{ex}
So sánh hai biểu thức: $67 - 5$ và $68 - 4$. Kết quả của biểu thức lớn hơn bằng bao nhiêu?
\shortans{64}
\loigiai{
$67 - 5 = 62$; $68 - 4 = 64$. Giá trị lớn hơn là $64$.
\textbf{Đáp số:} $64$.
}
\end{ex}

\begin{ex}
Bạn Hùng có 46 que tính, dùng 4 que xếp hình vuông. Hỏi Hùng còn lại bao nhiêu que tính?
\shortans{42}
\loigiai{
Số que tính Hùng còn lại là: $46 - 4 = 42$ (que tính).
\textbf{Đáp số:} $42$.
}
\end{ex}

\begin{ex}
Một bể nước có 89 lít nước. Người ta dùng hết 8 lít nước. Hỏi trong bể còn lại bao nhiêu lít nước?
\shortans{81}
\loigiai{
Số lít nước còn lại là: $89 - 8 = 81$ (lít).
\textbf{Đáp số:} $81$.
}
\end{ex}

\begin{ex}
Tính: $98 - 1 - 2 - 3$.
\shortans{92}
\loigiai{
Thực hiện: $98 - 1 = 97$; $97 - 2 = 95$; $95 - 3 = 92$.
\textbf{Đáp số:} $92$.
}
\end{ex}
